% GENERATED FILE. Edit canonical lessons, then run npm run generate:books.
% canonical-source-hash: fnv1a64:46322c0a72d37a5b
% canonical-lessons: MR-C466-samadhan, MR-C466-svasthya, MR-C466-nivant, MR-C466-visavne, MR-C466-khulne

\chapter{Peace of Mind}
\label{ch:mr-peace-of-mind}
\hlchaptermodality{466}

\begin{chapteropening}
\textbf{By the end of this chapter:} \emph{I can talk about satisfaction, ease and resting a while.}

Complete the last lesson of chapter 466: I can talk about satisfaction, ease and resting a while.
\end{chapteropening}

\section[\texorpdfstring{\mr{समाधान} (samādhān)}{समाधान (samādhān)}]{\mr{समाधान} (samādhān) — satisfaction; contentment}

\label{lesson:MR-C466-samadhan}



\begin{quote}
Before the new one: say the Marathi for to be baffled, then the Marathi for to be amazed.
\end{quote}

\subsection*{You'll want to know: \mr{समाधान}}
\textbf{\mr{समाधान}} — \emph{samādhān} — \textquotedblleft{}satisfaction; contentment\textquotedblright{}.

\begin{quote}
\mr{काम} \mr{पूर्ण} \mr{झाल्यावर} \mr{त्याला} \mr{समाधान} \mr{वाटलं}.
\end{quote}

\begin{quote}
\emph{kām pūrṇa jhālyāvar tyālā samādhān vāṭlaṁ.}
\end{quote}

\begin{quote}
He felt satisfied when the work was done.
\end{quote}

Neuter.

\mr{समाधान} \mr{वाटणे}, with the dative; \mr{समाधानी}, content.

\begin{cousinweb}
\textbf{\mr{समाधान}} (\emph{samādhān}) is a tatsama, taken from Sanskrit \emph{samādhāna}, \textquotedblleft{}putting together; reconciliation\textquotedblright{} (CDIAL 13210); in Marathi it is contentment.
\end{cousinweb}

\begin{grammarlens}[title={What you've built}]
One more word for describing what you lived and felt.
\end{grammarlens}

\subsection*{Your turn}
\begin{itemize}
\raggedright
  \item \emph{Say it:} \emph{samādhān}
  \item \emph{Say it:} \emph{samādhān}, once more
  \item \emph{Say it:} \emph{kām pūrṇa jhālyāvar tyālā samādhān vāṭlaṁ.} — He felt satisfied when the work was done.
  \item \emph{Recall:} say the Marathi for to be baffled, then the Marathi for to be amazed, then say \emph{samādhān} again
\end{itemize}

\begin{tcolorbox}[breakable,colback=teal!4,colframe=teal!35!black,title={Before you move on}]
What does \mr{समाधान} mean? (\textquotedblleft{}Satisfaction; contentment\textquotedblright{}.) Say it once more.
\end{tcolorbox}
\section[\texorpdfstring{\mr{स्वास्थ्य} (svāsthya)}{स्वास्थ्य (svāsthya)}]{\mr{स्वास्थ्य} (svāsthya) — well-being; ease}

\label{lesson:MR-C466-svasthya}



\begin{quote}
Before the new one: say the Marathi for to be amazed, then the Marathi for satisfaction.
\end{quote}

\subsection*{You'll want to know: \mr{स्वास्थ्य}}
\textbf{\mr{स्वास्थ्य}} — \emph{svāsthya} — \textquotedblleft{}well-being; ease\textquotedblright{}.

\begin{quote}
\mr{गोंगाटामुळे} \mr{मनाचं} \mr{स्वास्थ्य} \mr{बिघडतं}.
\end{quote}

\begin{quote}
\emph{goṅgāṭāmuḷe manācaṁ svāsthya bighaḍtaṁ.}
\end{quote}

\begin{quote}
Noise upsets one's peace of mind.
\end{quote}

Neuter.

\mr{मनाचं} \mr{स्वास्थ्य}, peace of mind; also health of body.

\begin{cousinweb}
\textbf{\mr{स्वास्थ्य}} (\emph{svāsthya}) is a tatsama, taken from Sanskrit \emph{svāsthya}, \textquotedblleft{}sound state of body or soul; ease, contentment\textquotedblright{} (Monier-Williams).
\end{cousinweb}

\begin{grammarlens}[title={What you've built}]
One more word for describing what you lived and felt.
\end{grammarlens}

\subsection*{Your turn}
\begin{itemize}
\raggedright
  \item \emph{Say it:} \emph{svāsthya}
  \item \emph{Say it:} \emph{svāsthya}, once more
  \item \emph{Say it:} \emph{goṅgāṭāmuḷe manācaṁ svāsthya bighaḍtaṁ.} — Noise upsets one's peace of mind.
  \item \emph{Recall:} say the Marathi for to be amazed, then the Marathi for satisfaction, then say \emph{svāsthya} again
\end{itemize}

\begin{tcolorbox}[breakable,colback=teal!4,colframe=teal!35!black,title={Before you move on}]
What does \mr{स्वास्थ्य} mean? (\textquotedblleft{}Well-being; ease\textquotedblright{}.) Say it once more.
\end{tcolorbox}
\section[\texorpdfstring{\mr{निवांत} (nivānt)}{निवांत (nivānt)}]{\mr{निवांत} (nivānt) — relaxed; at leisure}

\label{lesson:MR-C466-nivant}



\begin{quote}
Before the new one: say the Marathi for satisfaction, then the Marathi for well-being.
\end{quote}

\subsection*{You'll want to know: \mr{निवांत}}
\textbf{\mr{निवांत}} — \emph{nivānt} — \textquotedblleft{}relaxed; at leisure\textquotedblright{}.

\begin{quote}
\mr{सुट्टीच्या} \mr{दिवशी} \mr{आम्ही} \mr{निवांत} \mr{बसलो}.
\end{quote}

\begin{quote}
\emph{suṭṭīcyā divśī āmhī nivānt baslo.}
\end{quote}

\begin{quote}
On the holiday we sat around at leisure.
\end{quote}

Unchanging; also an adverb: \mr{निवांत} \mr{झोपणे}, to sleep undisturbed.

\begin{cousinweb}
\textbf{\mr{निवांत}} (\emph{nivānt}) is a Marathi word whose origin is not settled.
\end{cousinweb}

\begin{grammarlens}[title={What you've built}]
One more word for describing what you lived and felt.
\end{grammarlens}

\subsection*{Your turn}
\begin{itemize}
\raggedright
  \item \emph{Say it:} \emph{nivānt}
  \item \emph{Say it:} \emph{nivānt}, once more
  \item \emph{Say it:} \emph{suṭṭīcyā divśī āmhī nivānt baslo.} — On the holiday we sat around at leisure.
  \item \emph{Recall:} say the Marathi for satisfaction, then the Marathi for well-being, then say \emph{nivānt} again
\end{itemize}

\begin{tcolorbox}[breakable,colback=teal!4,colframe=teal!35!black,title={Before you move on}]
What does \mr{निवांत} mean? (\textquotedblleft{}Relaxed; at leisure\textquotedblright{}.) Say it once more.
\end{tcolorbox}
\section[\texorpdfstring{\mr{विसावणे} (visāvṇe)}{विसावणे (visāvṇe)}]{\mr{विसावणे} (visāvṇe) — to rest a while}

\label{lesson:MR-C466-visavne}



\begin{quote}
Before the new one: say the Marathi for well-being, then the Marathi for relaxed.
\end{quote}

\subsection*{You'll want to know: \mr{विसावणे}}
\textbf{\mr{विसावणे}} — \emph{visāvṇe} — \textquotedblleft{}to rest a while\textquotedblright{}.

\begin{quote}
\mr{चालून} \mr{थकल्यावर} \mr{आम्ही} \mr{झाडाखाली} \mr{विसावलो}.
\end{quote}

\begin{quote}
\emph{cālūn thaklyāvar āmhī jhāḍākhālī visāvlo.}
\end{quote}

\begin{quote}
Tired from walking, we rested a while under a tree.
\end{quote}

Past (he …): \mr{तो} \mr{विसावला}, \emph{to visāvlā}.

A short rest on the way: \mr{विसावा}, a halt.

\begin{cousinweb}
\textbf{\mr{विसावणे}} (\emph{visāvṇe}) is a tadbhava: CDIAL lists Marathi \emph{visavṇẽ}, \textquotedblleft{}to come to a halt, rest\textquotedblright{}, under Sanskrit \emph{viśramate}, \textquotedblleft{}rests\textquotedblright{} (CDIAL 11949).
\end{cousinweb}

\begin{grammarlens}[title={What you've built}]
One more word for describing what you lived and felt.
\end{grammarlens}

\subsection*{Your turn}
\begin{itemize}
\raggedright
  \item \emph{Say it:} \emph{visāvṇe}
  \item \emph{Say it:} \emph{visāvṇe}, once more
  \item \emph{Say it:} \emph{cālūn thaklyāvar āmhī jhāḍākhālī visāvlo.} — Tired from walking, we rested a while under a tree.
  \item \emph{Recall:} say the Marathi for well-being, then the Marathi for relaxed, then say \emph{visāvṇe} again
\end{itemize}

\begin{tcolorbox}[breakable,colback=teal!4,colframe=teal!35!black,title={Before you move on}]
What does \mr{विसावणे} mean? (\textquotedblleft{}To rest a while\textquotedblright{}.) Say it once more.
\end{tcolorbox}
\section[\texorpdfstring{\mr{खुलणे} (khulṇe)}{खुलणे (khulṇe)}]{\mr{खुलणे} (khulṇe) — to brighten up}

\label{lesson:MR-C466-khulne}



\begin{quote}
Before the new one: say the Marathi for relaxed, then the Marathi for to rest a while.
\end{quote}

\subsection*{You'll want to know: \mr{खुलणे}}
\textbf{\mr{खुलणे}} — \emph{khulṇe} — \textquotedblleft{}to brighten up\textquotedblright{}.

\begin{quote}
\mr{आजोबांना} \mr{पाहून} \mr{बाळाचा} \mr{चेहरा} \mr{खुलला}.
\end{quote}

\begin{quote}
\emph{ājobānnā pāhūn bāḷācā cehrā khullā.}
\end{quote}

\begin{quote}
The baby's face brightened at the sight of Grandfather.
\end{quote}

Past (he …): \mr{तो} \mr{खुलला}, \emph{to khullā}.

Of a face, a mood or a colour: \mr{रंग} \mr{खुलणे}, a colour sets something off.

\begin{cousinweb}
\textbf{\mr{खुलणे}} (\emph{khulṇe}) is inherited: CDIAL lists Marathi \emph{khulṇẽ}, \textquotedblleft{}to be open\textquotedblright{}, under the reconstructed \emph{*khull-} (CDIAL 3945); a face that opens up brightens.
\end{cousinweb}

\begin{grammarlens}[title={What you've built}]
One more word for describing what you lived and felt.
\end{grammarlens}

\subsection*{Your turn}
\begin{itemize}
\raggedright
  \item \emph{Say it:} \emph{khulṇe}
  \item \emph{Say it:} \emph{khulṇe}, once more
  \item \emph{Say it:} \emph{ājobānnā pāhūn bāḷācā cehrā khullā.} — The baby's face brightened at the sight of Grandfather.
  \item \emph{Recall:} say the Marathi for relaxed, then the Marathi for to rest a while, then say \emph{khulṇe} again
\end{itemize}

\begin{tcolorbox}[breakable,colback=teal!4,colframe=teal!35!black,title={Before you move on}]
What does \mr{खुलणे} mean? (\textquotedblleft{}To brighten up\textquotedblright{}.) Say it once more.
\end{tcolorbox}
